% généré par scripts/06_textes.py
\dosmeta{Année de classement : 380 (formule transmise par l’EDB (lue)). Références : CLE 1619 = ICUR I 3194 = ILCV 1464. Rattachement selon l’EDB : \emph{Incertae originis} (le volume I des ICVR réunit les inscriptions dont la provenance est perdue). Conservation selon l’EDB : \emph{villa Albani}. Support : Plaque de marbre. 69 × 128 cm (h. × l.). Inscription gravée. Reproductions signalées par l’EDB : reproduction, Perret 1851 (droits réservés ; non reproduites). Fiche EDB : Filippo Piazzolla, première version 2016. Notice du catalogue p.~\pageref{not:214}.}
\txtbloc{Texte transmis (EDCS, état 2022)}{\lnum{1}Hic quiescit ancilla dei qu(a)e de\par
\lnum{2}sua omnia possedit domum ista\par
\lnum{3}qu⸢a⸣m amic(a)e deflen(t) solaciumq(ue) requirunt\par
\lnum{4}pro hunc unum ora subolem qu⸢a⸣m superis\par
\lnum{5}titem re(li)quisti}
\txtbloc{Texte de l’EDB (EDB33145)}{\lnum{1}hic quescit ancilla dei ⌜q⌝ue de\par
\lnum{2}sua omnia possedit domum ista\par
\lnum{3}quem amice defien solaciumq(ue) requirun̂t\par
\lnum{4}pro hunc unum ora subolem quem superis\par
\lnum{5}titem requisti eterna requiem felicita/⟨ti⟩s cuasa man(e)bis IIIIX k⌜a⌝lendas octobris\par
\lnum{6}Cucurbitinus et Abu⌜nd⌝antius hic ⌜s⌝imul quiescit\par
\lnum{7}dd(ominis) nn(ostris) Gratiano V et Teodosio aaugg(ustis) ((aves coronam tenentes))}
